\section{Conclusion}

Public and private agent networks are two bundles of edge properties, not two values of one
variable. An edge is a triple---who may form it, what it carries, what is asserted about
it---and Twitter and WhatsApp are two corners of that space, not its two poles. Comparing
the corners moves all three axes at once, which is why the comparison has never told anyone
which component was doing the work, and why it returns $0.181$ or $0.314$ on the same
episodes depending on which pairing an author happens to call public and private.

Separated on one kernel, one runtime and one model, the three weigh $0.248$, $0.066$ and $0$
on a hidden-profile task family. The formation advantage is reach: it vanishes when the
roster already holds what a task needs. The semantics advantage is real, small, and charged
at 19k tokens a task. The attribute axis does not move, and a length-matched placebo of two
neutral lines is indistinguishable from an assertion of prior trust. \textbf{The first
result of a decomposition is that the space is smaller than it looks}: an axis whose effect
cannot be told apart from two lines of filler is not a dimension.

\textbf{On this task family the bounded network is weakly dominated, and we do not want to
be read as splitting the difference.} It never scores higher. It costs $\sim$19k more tokens
a task at every $E$ we ran---including $E=0$, where its roster holds everything the task
needs and it merely ties. It is slower. It fabricates values 58 times where the open arms
fabricated once, because the filter that keeps bad claims out also removes the second source
that catches them. Its single better column, disclosing to fewer counterparts, appears only
where it has already failed to reach anyone, which is not a privacy property. Against that,
one brokered hop recovers close to half the gap while every assertion \emph{about} trust
stays flat: an introduction changes the graph, an endorsement does not.

So if a private agent network is worth building, our measurements do not contain the reason,
and the reason is not information quality. Confidentiality, recourse over repeated
encounters, and resistance to counterparts with an incentive to lie are the three
candidates, and all three sit outside our scorer---each naming an experiment rather than a
hedge. Two further questions are the framework's own: whether the space is two-dimensional
everywhere or only on a task family that asks for reach, and whether an edge's position is
fixed at all. Open discovery buys the second source, a sandboxed first exchange makes a
stranger affordable, repeated success is the only thing that could make an attribute mean
anything, and shared state is worth its price only once an edge expects to be reused. If
contact sets narrow across repeats, a private network is what a public one becomes, and the
regimes were never alternatives. For an agentic web the implication is narrow and, we think,
useful: changing what a network permits changes what the collective does, and describing the
relationship does not---not at the layer where descriptions currently live.
